\documentclass[12pt,a4paper]{article}
\usepackage[utf8]{inputenc}
\usepackage[T1]{fontenc}
\usepackage{lmodern}
\usepackage[french]{babel}
\usepackage[a4paper,margin=2.3cm]{geometry}
\usepackage{amsmath,amssymb}
\usepackage{enumitem}
\setlist[itemize]{label=\textbullet,nosep}
\setlength{\parindent}{0pt}
\setlength{\parskip}{5pt}
\pagestyle{empty}
\newcommand{\question}[1]{\par\medskip{\large\bfseries #1}\par\smallskip}
\begin{document}

\begin{center}{\Large\bfseries Exercice de maison}\end{center}

On considère la suite $(U_n)$ définie par $U_0 = 3$ et, pour tout entier naturel $n$,
$U_{n+1} = U_n - 7$.

\question{1. Calcul de $U_1$ et $U_2$}
On utilise la relation $U_{n+1} = U_n - 7$.

Pour calculer $U_1$, on remplace $n$ par $0$ : \quad $U_1 = U_0 - 7 = 3 - 7 = -4$.

Pour calculer $U_2$, on remplace $n$ par $1$ : \quad $U_2 = U_1 - 7 = -4 - 7 = -11$.

Donc $\boxed{U_1 = -4 \ \text{ et } \ U_2 = -11}$.

\question{2. Montrons que $(U_n)$ est une suite arithmétique}
Pour tout entier naturel $n$, on a :
\[
U_{n+1} - U_n = (U_n - 7) - U_n = -7.
\]
La différence entre deux termes consécutifs est constante et égale à $-7$.
D'après la définition, $(U_n)$ est donc une suite arithmétique.

Donc $(U_n)$ est la suite arithmétique de \textbf{raison} $\boxed{r = -7}$ et de
\textbf{premier terme} $\boxed{U_0 = 3}$.

\question{3. Expression de $U_n$ en fonction de $n$}
$(U_n)$ est une suite arithmétique de premier terme $U_0 = 3$ et de raison $r = -7$.
Donc, pour tout entier naturel $n$ :
\[
U_n = U_0 + n\,r = 3 + n \times (-7).
\]
Donc, pour tout entier naturel $n$ : \quad $\boxed{U_n = 3 - 7n}$.

\question{4. Sens de variation de $(U_n)$}
$(U_n)$ est une suite arithmétique de raison $r = -7$.

Comme $r = -7$ et $-7 < 0$, alors la suite $(U_n)$ est \textbf{décroissante}.

\question{5. Calcul de $S = U_0 + U_1 + \dots + U_{100}$}
$S$ est la somme de termes consécutifs de la suite arithmétique $(U_n)$.
\begin{itemize}
  \item Nombre de termes : $100 - 0 + 1 = 101$.
  \item Premier terme : $U_0 = 3$.
  \item Dernier terme : $U_{100} = 3 - 7 \times 100 = 3 - 700 = -697$.
\end{itemize}
On a donc :
\[
S = (100 + 1) \times \frac{U_0 + U_{100}}{2}
  = 101 \times \frac{3 + (-697)}{2}
  = 101 \times \frac{-694}{2}
  = 101 \times (-347).
\]
Donc $\boxed{S = -35\,047}$.

\end{document}
